\section{Mechanisms of the arithmetic formulation}\label{sec:mechanisms}
%%% ====================================================================

\subsection{Prime decomposition and local Euler factors}\label{sec:primes}
%%% ====================================================================

Since $\Delta = 5$ is the only prime dividing the discriminant of
$\mathbb{Q}(\sqrt{5})$, the decomposition of primes in $\mathbb{Z}[\golden]$
is determined by the Legendre/Kronecker symbol $\chi_5(p) = (5/p)$:

\begin{proposition}[Local factors]\label{prop:local-factors}
The Dedekind zeta function of $\mathbb{Q}(\sqrt{5})$ has the Euler product
\begin{equation}\label{eq:euler-Qsqrt5}
  \zeta_{\mathbb{Q}(\sqrt{5})}(s)
  = \prod_{p\text{ prime}} f_p(s),
\end{equation}
where
\[
  f_p(s) = \begin{cases}
    (1-p^{-s})^{-1}   & \chi_5(p)=0  \text{ (ramified; only }p=5\text{)},\\
    (1-p^{-s})^{-2}   & \chi_5(p)=+1 \text{ (split; e.g.\ }p=11,19,29,\ldots\text{)},\\
    (1-p^{-2s})^{-1}  & \chi_5(p)=-1 \text{ (inert; e.g.\ }p=2,3,7,13,\ldots\text{)}.
  \end{cases}
\]
In particular, $p=2$ is \emph{inert} (Kronecker $(5/2)=-1$), so the factor
$\tfrac{1}{2}$ in $\Rpcf$ does not introduce a ramified place.

Each factor $f_p(s)$ is completely determined by the Frobenius lift
$\psi_p(\golden) = \golden^p$.
\end{proposition}

\begin{proof}
The classical theory of Dedekind zeta functions for quadratic fields; see
e.g.~\cite{Neukirch}.  The claim that $f_p$ is determined by $\psi_p$ follows
from the fact that $\chi_5(p) \equiv F_p \pmod{p}$
(\cref{prop:fib-split}): whether $\golden^p$ is
congruent to $\golden$ or $\bar\golden$ modulo $p$ determines the splitting.
\end{proof}

%%% ====================================================================
\subsection{Dirichlet factorization}\label{sec:factorisation}
%%% ====================================================================

\begin{theorem}[Factorisation]\label{thm:factorisation}
\begin{equation}\label{eq:factorisation}
  \zeta(s) = \frac{\zeta_{\mathbb{Q}(\sqrt{5})}(s)}{L(s,\chi_5)},
\end{equation}
where $L(s,\chi_5) = \prod_p(1-\chi_5(p)p^{-s})^{-1}$ is the Dirichlet
$L$-function of the symbol $\chi_5 = (5/\cdot)$.
\end{theorem}

\begin{proof}
The Dedekind zeta of a quadratic field $K$ factorises as
$\zeta_K(s)=\zeta(s)L(s,\chi_K)$, where $\chi_K$ is the Kronecker character
associated with $K/\mathbb{Q}$; for $K=\mathbb{Q}(\sqrt{5})$ this is~$\chi_5$.
See~\cite{Neukirch}, Ch.~VII.
\end{proof}

\begin{remark}[The Euler product determines $\zeta$ everywhere]%
  \label{rmk:euler-determines}
The identity $\prod_p(1-p^{-s})^{-1} = \sum_{n=1}^{\infty}n^{-s}$
for $\mathrm{Re}(s)>1$ (Euler~\cite{Euler}) follows from unique
factorisation of $\mathbb{N}$
(\cref{rmk:standard-id}).  The Dirichlet coefficients
$a_n = 1$ determine $\zeta(s)$ on all of $\mathbb{C}$ by analytic
continuation (identity theorem for Dirichlet series; Riemann~\cite{Riemann}).
Therefore the Euler product---and hence $G$---determines $\zeta(s)$
everywhere, including the critical strip $0<\mathrm{Re}(s)<1$ and
its zeros.
\end{remark}

\begin{remark}[Naturality of the factorisation]\label{rmk:natural-factorisation}
The factorisation $\zeta(s) = \zeta_{\mathbb{Q}(\sqrt{5})}(s)/L(s,\chi_5)$ is not specific
to $\mathbb{Q}(\sqrt{5})$: it is the $K=\mathbb{Q}(\sqrt{5})$
component of a natural isomorphism $\eta$ across all quadratic
extensions classified by $G$
(\cref{thm:dedekind-natural}).
The four components of $\eta$ correspond to the four fibers of the
Galois functor.
\end{remark}

%%% ====================================================================
\subsection{The fundamental formula}\label{sec:fundamental}
%%% ====================================================================

\begin{theorem}[Fundamental formula]\label{thm:L1}
\begin{equation}\label{eq:L1}
  L(1,\chi_5) = \frac{2\log\golden}{\sqrt{5}} = \frac{T}{\pi\sqrt{5}},
\end{equation}
where $T = 2\pi\log\golden$ is the period of $T^2_{\mathrm{PCF}}$.
\end{theorem}

\begin{proof}
This is the Dirichlet class-number formula for $\mathbb{Q}(\sqrt{5})$.
With class number $h=1$, regulator $R=\log\golden$, number of roots of unity
$w=2$, and discriminant $\Delta=5$:
\[
  \lim_{s\to 1}(s-1)\,\zeta_{\mathbb{Q}(\sqrt{5})}(s)
  = \frac{2^{r_1}(2\pi)^{r_2}hR}{w\sqrt{\Delta}}
  = \frac{4\cdot 1\cdot 1\cdot\log\golden}{2\sqrt{5}}
  = \frac{2\log\golden}{\sqrt{5}}.
\]
Since $\zeta_{\mathbb{Q}(\sqrt{5})}(s)=\zeta(s)L(s,\chi_5)$ and
$\lim_{s\to 1}(s-1)\zeta(s)=1$, we obtain $L(1,\chi_5)=2\log\golden/\sqrt{5}$.
The identity $T/(\pi\sqrt{5}) = 2\log\golden/\sqrt{5}$ follows from
$T=2\pi\log\golden$.
\end{proof}

\begin{remark}\label{rmk:pi-origin}
The \cref{eq:L1} exposes the role of the circle.
Writing $L(1,\chi_5) = T/(\pi\sqrt{5})$, we see that $\pi$ from the torus
period $T=2\pi\log\golden$ cancels against $\pi$ in the denominator,
leaving $2\log\golden/\sqrt{5}$.  The $\pi$---and hence the circle of
$T^2_{\mathrm{PCF}}$---is present from the outset.
\end{remark}

%%% ====================================================================
\subsection{Even zeta values}\label{sec:even}
%%% ====================================================================

\begin{theorem}[Even formula]\label{thm:even}
For $k \geq 1$:
\begin{equation}\label{eq:Lchi-even}
  L(2k,\chi_5)
  = \frac{(-1)^{k+1}\sqrt{5}\,(2\pi)^{2k}\,B_{2k,\chi_5}}
         {2\,(2k)!\,5^{2k}},
\end{equation}
where $B_{2k,\chi_5} \in \mathbb{Q}$ are the generalised Bernoulli numbers:
\[
  B_{n,\chi_5} = 5^{n-1}\sum_{a=1}^{5}\chi_5(a)\,B_n\!\left(\tfrac{a}{5}\right).
\]
Computed values: $B_{2,\chi_5}=\tfrac{4}{5}$,
$B_{4,\chi_5}=-8$, $B_{6,\chi_5}=\tfrac{804}{5}$, $B_{8,\chi_5}=-5776$.
\end{theorem}

\begin{proof}
Classical result: generalized Bernoulli formula for primitive even characters; see
e.g.~\cite{Washington}, Thm.~4.2.  The factor $\sqrt{5}$ is the Gauss sum
$\tau(\chi_5)=\sqrt{5}$.
\end{proof}

%%% ====================================================================
\subsection{The pentagonal source}\label{sec:pentagon}
%%% ====================================================================

The connection between $\golden$ and the odd $L$-values runs through
the pentagon.  We establish three identities---trigonometric, logarithmic,
and analytic---that collectively show how $\golden$ and $\pi$ organise
$L(s,\chi_5)$ for \emph{all} $s \geq 1$ via the same pentagonal structure.

\begin{proposition}[Pentagon identities]\label{prop:pentagon}
The golden ratio $\golden$ identifies the fundamental scales of the pentagon:
\begin{equation}\label{eq:pentagon-phi}
  2\cos\!\tfrac{\pi}{5} = \golden,
  \qquad
  2\cos\!\tfrac{2\pi}{5} = \golden^{-1}.
\end{equation}
These primary values extend symmetrically to the full set of four coordinates:
\begin{equation}\label{eq:pentagon-coords}
  2\cos\!\tfrac{\pi}{5} = \golden,\quad
  2\cos\!\tfrac{2\pi}{5} = \golden^{-1},\quad
  2\cos\!\tfrac{3\pi}{5} = -\golden^{-1},\quad
  2\cos\!\tfrac{4\pi}{5} = -\golden.
\end{equation}
\end{proposition}

\begin{proof}
The regular pentagon inscribed in the unit circle has diagonal-to-side
ratio $\golden$ (Euclid, \cite[Prop.~XIII.8]{heath1956thirteen}).  The half-angle of the central
subtended by one side is $\pi/5$, giving $2\cos(\pi/5)=\golden$.
The second identity follows from $\golden^{-1}=\golden-1$ and the
double-angle formula. Verified numerically to machine precision (see \cref{sec:verification}), yielding a zero residual.
\end{proof}

\begin{proposition}[Gauss logarithmic formula]\label{prop:gauss-log}
\begin{equation}\label{eq:gauss-log}
  \sum_{a=1}^{4}\chi_5(a)\,\log\!\Bigl(2\sin\tfrac{\pi a}{5}\Bigr)
  = -2\log\golden.
\end{equation}
\end{proposition}

\begin{proof}
Direct computation using $\chi_5(1)=\chi_5(4)=+1$,
$\chi_5(2)=\chi_5(3)=-1$, and the identities
$\sin(\pi/5)\sin(4\pi/5)=\sin^2(\pi/5)$,
$\sin(2\pi/5)\sin(3\pi/5)=\sin^2(2\pi/5)$.
The result follows from the product formula
$\prod_{a=1}^{4}\sin(\pi a/5) = \sqrt{5}/4$ and
$2\sin(\pi/5)\cdot 2\sin(2\pi/5) = \sqrt{5}$,
combined with \cref{prop:pentagon}.
Verified numerically (see \cref{sec:verification}).
\end{proof}

\begin{corollary}[Derivation of $L(1,\chi_5)$ from pentagon]\label{cor:L1}
The class-number formula (\cref{thm:L1}) follows from
\cref{prop:gauss-log} via the Gauss--Lerch formula for
primitive even characters with conductor $f$ and Gauss sum
$\tau(\chi_5)=\sqrt{5}$:
\begin{equation}\label{eq:gauss-lerch}
  L(1,\chi_5)
  = -\frac{\tau(\chi_5)}{f}\sum_{a=1}^{f-1}\chi_5(a)
      \log\!\Bigl(2\sin\tfrac{\pi a}{f}\Bigr)
  = -\frac{\sqrt{5}}{5}\cdot(-2\log\golden)
  = \frac{2\log\golden}{\sqrt{5}}
  = \frac{T}{\pi\sqrt{5}}.
\end{equation}
The pentagon generates $\golden$ (\cref{prop:pentagon}),
which enters $L(1,\chi_5)$ through the Gauss sum over
$\sin(\pi a/5)$.
\end{corollary}

\begin{proposition}[Digamma formula]\label{prop:digamma}
Let $\psi = \Gamma'/\Gamma$ denote the digamma function.  Then
\begin{equation}\label{eq:digamma-sum}
  \sum_{a=1}^{4}\chi_5(a)\,\psi\!\bigl(\tfrac{a}{5}\bigr)
  = -2\sqrt{5}\,\log\golden,
\end{equation}
and consequently $L(1,\chi_5) = -\tfrac{1}{5}\sum_{a=1}^{4}
\chi_5(a)\,\psi(a/5) = 2\log\golden/\sqrt{5}$.
\end{proposition}

\begin{proof}
The digamma formula at rational arguments gives
$\psi(a/5) = -\gamma - \log(5) - \frac{\pi}{2}\cot(\pi a/5)
+ 2\sum_{k=1}^{2}\cos(2\pi ka/5)\log\sin(\pi k/5)$
(Gauss's digamma theorem).  The $\chi_5$-weighted sum eliminates
$\gamma$ and $\log(5)$ (since $\sum\chi_5(a)=0$) and reduces to the
logarithmic formula of \cref{prop:gauss-log}.
Verified numerically (see \cref{sec:verification}).
\end{proof}


\begin{proposition}[Pentagon dictates $\chi_5$]\label{prop:pentagon-dictates}
The Artin character $\chi_5$ of $\mathbb{Q}(\sqrt{5})/\mathbb{Q}$ is
completely determined by the pentagon through:
\begin{equation}\label{eq:chi5-from-pentagon}
  \chi_5(a)
  = \frac{\log|2\cos(\pi a/5)|}{\log\golden}
  \in \{+1,-1\},
  \qquad a = 1,2,3,4.
\end{equation}
Equivalently:
\begin{equation}\label{eq:norm-cosine-phi}
  |2\cos(\pi a/5)| = \golden^{\chi_5(a)}.
\end{equation}
The four pentagonal cosines $2\cos(\pi a/5)\in\{\pm\golden,\pm\golden^{-1}\}$
are the four images of the Galois functor
$G\colon\Ztwenty\to\mathbb{Z}_2\times\mathbb{Z}_2$
evaluated at the residues $\{1,2,3,4\}\subset\mathbb{Z}/5\mathbb{Z}$.
\end{proposition}

\begin{proof}
\Cref{prop:pentagon} gives:
$2\cos(\pi/5)=\golden$ and $2\cos(2\pi/5)=\golden^{-1}$.
By symmetry of cosine: $2\cos(3\pi/5)=-\golden^{-1}$ and $2\cos(4\pi/5)=-\golden$.
Hence $|2\cos(a\pi/5)|\in\{\golden,\golden^{-1}\}$ for all $a\in\{1,2,3,4\}$.
Taking logarithms: $\log|2\cos(a\pi/5)|\in\{\log\golden,-\log\golden\}$.
Dividing by $\log\golden$: $\log|2\cos(a\pi/5)|/\log\golden\in\{+1,-1\}$.

The identification with $\chi_5(a)$: for $a\in\{1,4\}$ (quadratic residues
mod~5), $|2\cos(a\pi/5)|=\golden>\golden^{-1}$, so the ratio is $+1=\chi_5(a)$.
For $a\in\{2,3\}$ (non-residues), $|2\cos(a\pi/5)|=\golden^{-1}<\golden$,
so the ratio is $-1=\chi_5(a)$.
Verified numerically (see \cref{sec:verification}).
\end{proof}

\begin{remark}[Connection to the Galois functor]\label{rmk:pentagon-G}
The pentagonal cosines of \cref{prop:pentagon-dictates}
reappear as the binary $(0,1)$-component of the Galois functor
$G\colon\Ztwenty\to\mathbb{Z}_2\times\mathbb{Z}_2$
(\cref{thm:bridge3} below): $\chi_5(a)$ is simultaneously
the pentagonal fingerprint~\eqref{eq:chi5-from-pentagon} and the
second coordinate of $G(a\bmod 20)$.
\end{remark}

\begin{remark}[The pentagonal dictate]
  \label{rmk:pentagonal-dictate}
\Cref{prop:pentagon-dictates} makes explicit what
``$\golden$ dictates $\chi_5$'' means concretely.
The chain is:
\[
\begin{split}
  &\golden = 2\cos\tfrac{\pi}{5}
  \;\Rightarrow\;
  \chi_5(a) = \tfrac{\log|2\cos(a\pi/5)|}{\log\golden}\\[4pt]
  &\quad\Rightarrow\;
  L(s,\chi_5) = \tfrac{1}{5^s}\textstyle\sum_a \chi_5(a)\,\zeta(s,\tfrac{a}{5})
  \;\Rightarrow\;
  \zeta(2k+1) = \tfrac{\zeta_{\mathbb{Q}(\sqrt{5})}(2k+1)}{L(2k+1,\chi_5)}.
\end{split}\]
The values $\zeta(2k+1)$ that appear as free parameters in the EFT of
Elvang et al.\ are not free: they are Hurwitz projections at the
pentagonal points $\{a/5\}$ with weights $\chi_5(a)$, which are
themselves the logarithmic fingerprints of $\golden$ in the pentagon.
The ``dictum'' is $\golden=2\cos(\pi/5)$; everything else follows.
\end{remark}

\begin{theorem}[Hurwitz representation and pentagonal source]%
  \label{thm:hurwitz}
For all $s$ with $\mathrm{Re}(s)>1$ (and by analytic continuation
for $s\neq 1$):
\begin{equation}\label{eq:hurwitz-rep}
  L(s,\chi_5)
  = \frac{1}{5^s}\sum_{a=1}^{4}\chi_5(a)\,\zeta(s,\tfrac{a}{5}),
\end{equation}
where $\zeta(s,x)=\sum_{n=0}^{\infty}(n+x)^{-s}$ is the Hurwitz
zeta function.  The four arguments $a/5$ ($a=1,2,3,4$) are the
\emph{pentagonal points}, and $\chi_5$ is their Artin weight.

The three kernels at special values:
\begin{align}
  s = 1\colon\quad
    & \zeta(s,a/5)\xrightarrow{s\to 1} -\psi(a/5) - \log 5
    \;\Rightarrow\; L(1,\chi_5) = \tfrac{2\log\golden}{\sqrt{5}},
    \label{eq:kernel-s1} \\
  s = 2k\colon\quad
    & \zeta(2k,a/5)\text{ via Bernoulli}
    \;\Rightarrow\; L(2k,\chi_5) = \sqrt{5}\cdot\pi^{2k}\cdot\mathbb{Q},
    \label{eq:kernel-even} \\
  s = 2k+1,\; k\geq 1\colon\quad
    & \zeta(2k+1,a/5) = \sum_{n=0}^{\infty}(n+a/5)^{-(2k+1)}.
    \label{eq:kernel-odd}
\end{align}
In all three cases the organising structure is identical: the same
$\chi_5(a)$-weights over the same pentagonal points $\{a/5\}$.
\end{theorem}
\begin{proof}
\Cref{eq:hurwitz-rep} is well-established (\cite[Thm.~12.6]{apostol1976introduction}).
The kernel at $s=1$ follows from $\zeta(s,x) = 1/(s-1) + \psi(x)+O(s-1)$
combined with $\sum\chi_5(a)=0$ (so the $1/(s-1)$ pole and $\log 5$
cancel), giving \cref{prop:digamma} and hence
\cref{cor:L1}.  The kernel at even $s$ gives
\cref{thm:even}.  The odd kernel~\cref{eq:kernel-odd} is the
direct Hurwitz series, convergent for $s>1$.
\end{proof}

\begin{theorem}[Pentagonal determination of odd zeta values]%
  \label{thm:pentagon-odd}
For $k\geq 0$, the factorisation $\zeta(2k+1) = \zeta_{\mathbb{Q}(\sqrt{5})}(2k+1)/L(2k+1,\chi_5)$
expresses $\zeta(2k+1)$ as the ratio of two $\chi_5$-projections of
Hurwitz zeta at pentagonal points.  All data are generated by:
\begin{enumerate}
\item $2\cos(\pi/5) = \golden$ \emph{(pentagonal geometry generates $\golden$)},
\item $\sum_{a=1}^{4}\chi_5(a)\log(2\sin\frac{\pi a}{5}) = -2\log\golden$
  \emph{(the logarithmic signature of $\golden$ in the pentagon)},
\item $\psi_p(\golden)=\golden^p \;\Rightarrow\; \chi_5(p) = $ Artin character
  \emph{(Frobenius connects $\golden$ to the L-function character)}.
\end{enumerate}
The case $k=0$ yields the exact closed form
$L(1,\chi_5) = T/(\pi\sqrt{5})$.
For $k\geq 1$, the determination is structural: $L(2k+1,\chi_5)$ and
$\zeta_{\mathbb{Q}(\sqrt{5})}(2k+1)$ are both projections of the same pentagonal Hurwitz
structure, organised by $\chi_5 = $ Artin character of
$\mathbb{Q}(\sqrt{5})/\mathbb{Q}$.  No additional parameter beyond
$\golden$, $\pi$, and the pentagonal arithmetic of $\mathbb{Q}(\sqrt{5})$
is required.
\end{theorem}
\begin{proof}
Items (1)--(3) are \cref{prop:pentagon,prop:gauss-log} and \cref{thm:bridge3} respectively.
The factorisation is \cref{thm:factorisation}.
The Hurwitz representation is \cref{thm:hurwitz}.
The $k=0$ closed form is \cref{cor:L1}.
For $k\geq 1$: both $\zeta_{\mathbb{Q}(\sqrt{5})}(2k+1) = \zeta(2k+1)\cdot L(2k+1,\chi_5)$
and $L(2k+1,\chi_5) = 5^{-(2k+1)}\sum\chi_5(a)\zeta(2k+1,a/5)$
are determined by the Hurwitz series at pentagonal points weighted
by $\chi_5$.  Neither factor has a simpler closed form (\cref{op:closed-form}), but their ratio $\zeta(2k+1)$ is
fixed by the arithmetic of $\Rpcf$ without any free parameter.
\end{proof}

\begin{remark}[Structural vs.\ closed-form determination]\label{rmk:structural}
\Cref{thm:pentagon-odd} establishes \emph{structural} determination: the
pentagonal lattice $\{a/5\}$, the Artin character $\chi_5$, and the
Hurwitz kernel $\zeta(2k+1,\cdot)$ are all fixed by $\golden$ and $\pi$
without additional parameters.  A closed-form expression
$\zeta(2k+1) = r_k\pi^{2k+1}$ for $r_k\in\mathbb{Q}$ would be stronger;
whether it exists for $k\geq 1$ is an open problem related to Apéry's
theorem and its generalisations.  What the present paper establishes is
that the \emph{source}---the organising structure---is the pentagon
generated by $\golden$ and $\pi$, not free arithmetic data.
\end{remark}

% ------------------------------------------------------------------
\subsubsection{Numerical case study: \texorpdfstring{$\zeta(3)$}{zeta(3)}}\label{ssec:zeta3-example}
% ------------------------------------------------------------------

We trace the full chain from the pentagon to
$\zeta(3) = 1.2020569\ldots$ for the first few primes.

\textbf{Step~1: Pentagon $\to$ character.}
$\golden = 2\cos(\pi/5)$.  The Artin character
$\chi_5(p) = (5/p)$ classifies primes:
$\chi_5(2)=-1$, $\chi_5(3)=-1$, $\chi_5(7)=-1$,
$\chi_5(11)=+1$, $\chi_5(13)=-1$, $\chi_5(19)=+1$.

\textbf{Step~2: Character $\to$ local factors at $s=3$.}
Split ($\chi_5=+1$): $f_p(3) = (1-p^{-3})^{-2}$.
Inert ($\chi_5=-1$): $f_p(3) = (1-p^{-6})^{-1}$.
Ramified ($p=5$): $f_5(3) = (1-5^{-3})^{-1}$.
Explicitly:
\[
  f_2(3) = \frac{64}{63},\quad
  f_3(3) = \frac{729}{728},\quad
  f_5(3) = \frac{125}{124},\quad
  f_{11}(3) = \left(\frac{1331}{1330}\right)^{\!2}.
\]

\textbf{Step~3: Euler product $\to$ $\zeta(3)$.}
$\zeta_{\mathbb{Q}(\sqrt{5})}(3) = \prod_p f_p(3)$ and
$L(3,\chi_5) = \prod_p(1-\chi_5(p)p^{-3})^{-1}$.
Their ratio gives
\[
  \zeta(3) = \frac{\zeta_{\mathbb{Q}(\sqrt{5})}(3)}{L(3,\chi_5)} = 1.2020569\ldots
\]
Every factor in both products is determined by $\chi_5(p)$,
which is determined by $F_p\bmod p$
(\cref{prop:fib-split}), which is determined
by $\golden^p = F_p\golden + F_{p-1}$
(\cref{def:frobenius}).
No parameter is chosen; the pentagon dictates the value.
Verified numerically (see \cref{sec:verification}).

The mechanism extends to all odd values.  \Cref{tab:zeta-values}
displays the decomposition $\zeta(2k+1) = \zeta_{\mathbb{Q}(\sqrt{5})}(2k+1)/L(2k+1,\chi_5)$
alongside the even values $\zeta(2k) = r_k\pi^{2k}$, illustrating the
two regimes: even values governed by $\pi$ alone, odd values by both
$\golden$ and $\pi$.

\begin{table}[ht]
\centering
\caption{Zeta values from $s = 2$ to $s = 13$.  Even values (shaded)
are rational multiples of $\pi^{2k}$; odd values are ratios
$\zeta_{\mathbb{Q}(\sqrt{5})}(s)/L(s,\chi_5)$, both factors determined by
$\chi_5$ and hence by $\golden$.  All entries verified at
50-digit precision (see \cref{sec:verification}).}
\label{tab:zeta-values}
\begin{tabular}{rcccc}
\toprule
$s$ & $\zeta_{\mathbb{Q}(\sqrt{5})}(s)$ & $L(s,\chi_5)$ & $\zeta(s)$ & Source \\
\midrule
\rowcolor[gray]{0.93}
 2 & --- & --- & $\pi^2/6 \approx 1.6449$ & $\pi$ \\
 3 & $1.02755$ & $0.85482$ & $1.20206$ & $\golden,\pi$ \\
\rowcolor[gray]{0.93}
 4 & --- & --- & $\pi^4/90 \approx 1.0823$ & $\pi$ \\
 5 & $1.00133$ & $0.96567$ & $1.03693$ & $\golden,\pi$ \\
\rowcolor[gray]{0.93}
 6 & --- & --- & $\pi^6/945 \approx 1.0173$ & $\pi$ \\
 7 & $1.00007$ & $0.99179$ & $1.00835$ & $\golden,\pi$ \\
\rowcolor[gray]{0.93}
 8 & --- & --- & $\pi^8/9450 \approx 1.0041$ & $\pi$ \\
 9 & $1.000004$ & $0.99800$ & $1.00201$ & $\golden,\pi$ \\
\rowcolor[gray]{0.93}
10 & --- & --- & $\pi^{10}/93555 \approx 1.00099$ & $\pi$ \\
11 & $1.0000003$ & $0.99951$ & $1.00049$ & $\golden,\pi$ \\
\rowcolor[gray]{0.93}
12 & --- & --- & $691\pi^{12}/638512875 \approx 1.00025$ & $\pi$ \\
13 & $1.00000002$ & $0.99988$ & $1.00012$ & $\golden,\pi$ \\
\bottomrule
\end{tabular}
\end{table}

%%% ====================================================================
\subsection{The period bridge}\label{sec:period}
%%% ====================================================================

The two structures of $T^2_{\mathrm{PCF}}$ interact precisely through
the period $T$:

\begin{center}
\begin{tabular}{lll}
\toprule
Structure & Object & Enters via \\
\midrule
Multiplicative (ring) & $\psi_p(\golden)=\golden^p$ & Euler product factors $f_p(s)$ \\
Additive/circular (torus) & $T = 2\pi\log\golden$ & $L(1,\chi_5) = T/(\pi\sqrt{5})$ \\
\bottomrule
\end{tabular}
\end{center}

The \cref{eq:L1} is the bridge: $\pi$ (circle)
and $\log\golden$ (Frobenius regulator) combine in $T$, then $\pi$ cancels
to leave $2\log\golden/\sqrt{5}$.  At $s=2k$ the $\pi$ reappears
explicitly via~\cref{eq:Lchi-even}.  At $s=2k+1$ for $k\geq 1$ the
same two structures produce $\zeta(2k+1)$ through their ratio.

The table displays two faces of $T^2_{\mathrm{PCF}}$: multiplicative
(Frobenius, Euler product) and additive (period, circle).
We now show that both arise from the self-measurement
$\golden^2 = \golden + 1$.

% ------------------------------------------------------------------
\subsubsection{Dimensional emergence}\label{ssec:dim-emergence}
% ------------------------------------------------------------------

\begin{proposition}[Dimensional chain from self-measurement]%
  \label{prop:dim-chain}
The algebraic relations $i^2 = -1$ and $\golden^2 = \golden + 1$
generate a dimensional chain
\begin{equation}\label{eq:dim-chain}
  \underbrace{\mathbb{R}}_{d=1}
  \;\xrightarrow{\;\times\, i\;}
  \underbrace{\mathbb{C}}_{d=2}
  \;\xrightarrow{\;\times\,\golden\;}
  \underbrace{E^3}_{d=3},
\end{equation}
where $E^3 = \{(x,\,y,\,\golden y) : x,y\in\mathbb{R}\}$
is the golden extension of\/ $\mathbb{C}$. Under the coupling $z = \golden y$, 
the characteristic equation $\golden^2 = \golden + 1$ manifests 
as the self-measurement identity:
\begin{equation}\label{eq:self-measurement}
  z^2 = z\cdot y + y^2.
\end{equation}
\end{proposition}

\begin{proof}
(i) \textbf{$\mathbb{R}\to\mathbb{C}$:} This is the foundational construction of\/ $\mathbb{C}$ as
$\mathbb{R}[x]/(x^2+1)$. Geometrically, $i^2 = -1$ signifies that squaring 
a 1D generator produces a result outside\/ $\mathbb{R}$; the discrepancy 
constitutes the imaginary axis.

(ii) \textbf{$\mathbb{C}\to E^3$:} The ambient space $\mathbb{R}^3$ has canonical basis
$\{e_1,e_2,e_3\}$. The coupling $z = \golden y$ identifies the
third coordinate with the second scaled by~$\golden$, giving
$E^3 = \{x\,e_1 + y(e_2 + \golden\,e_3) : x,y\in\mathbb{R}\}$.
This is a two-dimensional subspace of~$\mathbb{R}^3$: the third
ambient dimension exists but is slaved to the second.
The irrationality of $\golden$ ensures the coupling is non-degenerate
(\cref{prop:phi-unique}): $\golden\notin\mathbb{Q}$ means $z = \golden y$ 
cannot be reduced to a rational relation between coordinates.

(iii) \textbf{Self-measurement:} Setting $z = \golden y$, direct computation 
using $\golden^2 = \golden + 1$ yields:
\[ z^2 = (\golden y)^2 = \golden^2 y^2 = (\golden+1)y^2 = \golden y\cdot y + y^2 = zy + y^2. \]
The square of the third coordinate (a 2D operation applied to 
the new dimension) folds back onto a linear combination of 
lower-dimensional terms $z y$ and $y^2$. This is the geometric 
content of $\golden^2 = \golden + 1$: the new dimension, 
when measured against itself, refers back to the dimensions 
that generated it.
\end{proof}

\begin{remark}[The Frobenius lift is arithmetic self-measurement]%
  \label{rmk:frobenius-selfmeas}
The self-measurement equation~\eqref{eq:self-measurement} has exactly
the structure of the Frobenius lift: setting $y=1$,
\[
  z^2 = z + 1
  \qquad\longleftrightarrow\qquad
  \golden^2 = \golden + 1,
\]
and for prime $p$:
\[
  \golden^p = F_p\,\golden + F_{p-1}.
\]
Each prime $p$ performs the same operation as the dimensional
extension---it raises $\golden$ to a higher power and the result
folds back onto the basis $\{\golden, 1\}$.  The splitting
character $\chi_5(p) \equiv F_p\pmod{p}$
(\cref{prop:fib-split}) records the outcome:
\begin{itemize}
\item $\chi_5(p) = +1$ (split): $\golden^p\equiv\golden\pmod{p}$.
  The prime sees $\golden$ as itself---the measurement is
  one-dimensional.
\item $\chi_5(p) = -1$ (inert): $\golden^p\equiv\bar\golden\pmod{p}$.
  The prime sees $\golden$ as its conjugate---the measurement
  detects the second dimension.
\end{itemize}
The Euler product $\zeta_{\mathbb{Q}(\sqrt{5})}(s) = \prod_p f_p(s)$
(\cref{prop:local-factors}) is the collective
measurement of all primes, converging to the Dedekind zeta
of $\mathbb{Q}(\sqrt{5})$---an inherently two-dimensional object.
\end{remark}

\begin{remark}[Dedekind factorisation as dimensional projection]%
  \label{rmk:dimensional-projection}
Under the identification of \cref{prop:dim-chain} and
\cref{rmk:frobenius-selfmeas}, the factorisation
\[
  \zeta(s) = \frac{\zeta_{\mathbb{Q}(\sqrt{5})}(s)}{L(s,\chi_5)}
\]
(\cref{thm:factorisation}) admits a reading as dimensional
projection: the one-dimensional $\zeta(s)$ (primes over
$\mathbb{Q}$) is the quotient of the two-dimensional $\zeta_{\mathbb{Q}(\sqrt{5})}(s)$
(primes over $\mathbb{Q}(\sqrt{5})$) by the character $L(s,\chi_5)$
that encodes how the second dimension projects onto the first.

In this reading, the even values
$\zeta(2k) = \pi^{2k}\cdot\mathbb{Q}$ (Euler) are pure
circle geometry---the second dimension already closed.
The odd values $\zeta(2k+1)$
(\cref{thm:pentagon-odd}) are the mixed terms where the
1D generator ($\golden$ via $\chi_5$) and the 2D geometry
($\pi$ via the pentagonal points $a/5$) couple.
They appear free in the EFT of
Elvang et al.~\cite{Elvang2026} because the EFT is a
four-dimensional quantum field theory that does not access
$T^2_{\mathrm{PCF}}$; the Veneziano amplitude resolves
them because it is a worldsheet object---inherently
two-dimensional.

The regulator $R = \log\golden$
(\cref{prop:regulator-unifier}) measures the
informational cost of this emergence: the period $T = 2\pi R$
is the circumference of the new dimension, the $L$-value
$L(1,\chi_5) = 2R/\sqrt{5}$ is the arithmetic size of the
2D field over the 1D field, and the entropy
$S_{\mathrm{BH}}/k_B = \lambda_{\log}\cdot R$ is the
information content of the projection.
Three formulas, one invariant, one geometric fact: the second
dimension costs $\log\golden$ to build.
\end{remark}


%%% ====================================================================
